
Let \(Y'\) be an open subset of \(Y\), and put
\[
  U'=q^{-1}(Y'),\qquad X'_0=p^{-1}(Y'),
\]
and so forth.\footnote{\label{note:sga3-V-33}Editor's note: The argument that follows has
been expanded, in particular the verification that \(U'\) is a
quasi-section of the groupoid induced on \(X'_0\).}
Then \(U'\) is an open subset of \(U\) saturated for the equivalence
relation defined by the groupoid \(U_\ast\). It follows from
\hyperref[lem:sga3-V-1.1]{Lemmas~1.1} and
\hyperref[sga3:exp5:lemma:1.2]{1.2} that \(Y'\) is the cokernel, both
in \((\mathbf{Sch}/S)\) and in the category of ringed spaces, of the
groupoid induced by \(U_\ast\) on \(U'\). Likewise, \(X'_0\) is an
open subset of \(X_0\) saturated for the equivalence relation defined
by \(X_\ast\), and we have the following commutative diagram, in which
both squares are Cartesian:
\SGAThreeDiagram[0.58\textwidth]
  {figures/exp5/p273_lemma61_restricted_groupoids_cartesian_corrected.png}
  {The restricted groupoids over \(X'_0\) and \(X_0\).}
  {fig:sga3-V-6.1-restricted-groupoids}
  {Expos\'e V, component-local p.~23 / combined-reader p.~285 / re-edited header p.~273; lower-right \(U\) corrected from source-PDF \(U'\)}
Here \(\widetilde d_1\) is surjective and finite locally free. On the
other hand, let \(x\in U'\). Since \(U\) is a quasi-section, the set
\[
  E=d_0d_1^{-1}(x)\cap U
\]
is finite and is contained in an affine open subset \(W\) of \(U\).
Then \(E'=E\cap U'\) is a finite set contained in the quasi-affine
open subset \(W\cap U'\). Consequently there is an affine open subset
\(W'\) of \(W\cap U'\) that contains \(E'\). This proves that \(U'\)
is a quasi-section of the groupoid \(X'_\ast\) induced by \(X_\ast\)
on \(X'_0\). The first assertion of \hyperref[lem:sga3-V-6.1-i]{part~(i)}, applied to \(X'_\ast\)
and \(U'\), now shows that \(Y'\) is the cokernel of \(X'_\ast\) in
\((\mathbf{Sch}/S)\).

In particular, for every \(S\)-scheme \(T\), the sequence
\begin{equation}
  T(Y')
  \xrightarrow{\,T(p|_{X'_0})\,}
  T(X'_0)
  \underset{T(d_0|_{X'_1})}
    {\overset{T(d_1|_{X'_1})}{\rightrightarrows}}
  T(X'_1)
  \label{eq:sga3-V-6.1-T-exact}
\end{equation}
is exact. If \(T\) is the affine line \(\mathbf G_{a,S}\)
(\SGARefLink{sga3:i:subsection:4-3}{Expos\'e I, 4.3}), this sequence identifies with
\begin{equation}
\begin{aligned}
  \Gamma(Y',\mathcal O_Y)
  &\longrightarrow
  \Gamma(p^{-1}(Y'),\mathcal O_{X_0})
  \underset{\delta_0}{\overset{\delta_1}{\rightrightarrows}}
  \Gamma(d_0^{-1}p^{-1}(Y'),\mathcal O_{X_1})\\
  &\hspace{9em}=
  \Gamma(d_1^{-1}p^{-1}(Y'),\mathcal O_{X_1}),
\end{aligned}
\label{eq:sga3-V-6.1-sections-exact}
\end{equation}
which is therefore exact for every open subset \(Y'\). This completes
the proof of Lemma~\hyperref[lem:sga3-V-6.1-i]{6.1(i)}.

\medskip
\phantomsection\label{par:sga3-V-6-d}%
\noindent\textbf{d) Proof of \hyperref[lem:sga3-V-6.1-iv]{(iv)}.}---%
If \((d_0,d_1)\) is an equivalence pair, then so is
\((u_0,u_1)\). It follows from
Theorem~\hyperref[thm:sga3-V-4.1]{4.1} that
\[
  u_0\boxtimes u_1:U_1\longrightarrow U\times_Y U
\]
is an isomorphism, and hence so is \(v_0\boxtimes v_1\); see the
Cartesian squares in part~\hyperref[par:sga3-V-6-b]{(b)}. Since
\(v\times v\) is faithfully flat and quasi-compact,
\(d_0\boxtimes d_1\) is an isomorphism (SGA~1, VIII, \SGARefLink{sga3:viii:corollary:5:4}{5.4}).

Moreover, if \(d_0\) is flat, then \(u\) is flat. The morphism \(q\)
is flat by Theorem~\hyperref[thm:sga3-V-4.1]{4.1}, so \(r\) is flat
as well. Since \(v\) is faithfully flat, \(p\) is flat; being also
surjective, it is faithfully flat.
